% =============================================================
%  ch07_zykloidgetriebe.tex  --  Zykloidgetriebe
% =============================================================

\chapter{Zykloidgetriebe}

Ein Zykloidgetriebe erzeugt auf kompaktem Bauraum eine hohe Untersetzung bei
geringem Spiel (Backlash) und hoher Stoßfestigkeit -- ideal als ``letzte
Stufe'' vor dem Robotergelenk. Eine exzentrisch gelagerte Zykloidscheibe mit
$n$ Nocken wälzt in einem Ring mit $N$ Bolzen ab. Die Untersetzung ist
\begin{equation}
i = \frac{N}{N - n}.
\end{equation}
Mit einer Nockendifferenz von eins ($N - n = 1$) folgt $i = N$: 11 Ringbolzen
und eine 10-Nocken-Scheibe ergeben 11:1.

Für die Simulation ist die innere Mechanik irrelevant -- das Getriebe wird als
\emph{ein} angetriebenes Gelenk mit Übersetzung modelliert (siehe
Teil~IV, ros2\_control-Transmission). Wichtig für die reale Auslegung: das
übertragbare Ausgangsmoment, die Steifigkeit und das Restspiel, die du später
in die Gelenklimits einträgst.
